\centerline{\bf \Huge Лемма о трезубце.}
\centerline{\bf \Huge }

\begin{flushright}
        \scriptsize{}
\end{flushright}


\noindent\textbf{Лемма о трезубце.}  Пусть у треугольника \( ABC \) точка \( I \) — центр вписанной окружности, точка \( I_a \) — центр вневписанной окружности, противоположной вершине \( A \), а точка \( L \) — точка пересечения отрезка \( II_a \) с дугой описанной окружности. Тогда точка \( L \) равноудалена от \( I \), \( I_a \), \( B \) и \( C \).

\noindent\textbf{Внешняя лемма о трезубце.} Пусть $I_B$ и $I_C$ --- центры вневписанных окружностей треугольника $ABC$, касающихся сторон $AC$ и $AB$ соответственно. Тогда точки $B, C, I_C, I_B$  лежат на одной окружности с центром в середине "большей" дуги $BC$ описанной окружности треугольника $ABC$.

\problem{1}
Вокруг прямоугольного треугольника $ABC$ с прямым углом $C$ описана окружность, на меньших дугах $AC$ и $BC$   взяты их середины --- $K$ и $P$  соответственно. Отрезок $KP$ пересекает катет $AC$ в точке $N$.   Центр вписанной окружности треугольника $ABC$  ---  $I$.   Найти угол $NIC$.

\problem{2}
В окружность вписан четырёхугольник $ABCD$ Отметили центры окружностей, впи-
санных в треугольники $BCD$, $CDA$, $DAB$, $ABC$. Докажите, что отмеченные точки являются вершинами прямоугольника.

\problem{3}
Дана равнобокая трапеция $ABCD$ с основаниями $BC$ и $AD$. В треугольники $ABC$ и $ABD$ вписаны окружности с центрами $I$ соответственно. Докажите, что прямая $I_1I_2$   перпендикулярна $BC$.

\problem{4}
$\bigl[\textbf{Финал ВсОШ, 2022, 9.2}\bigr]$
Биссектрисы треугольника $ABC$ пересекаются в точке $I$, внешние биссектрисы его углов $B$ и $C$ пересекаются в точке $J$. Окружность $\omega_b$ с центром в точке $O_b$ проходит через точку $B$ и касается прямой $CI$ в точке $I$. Окружность $omega_c$ с центром в точке $O_c$ проходит через точку $C$ и касается прямой $BI$ в точке $I$. Отрезки $O_bO_c$ и $IJ$ пересекаются в точке $K$. Найдите отношение $IK : KJ$.